\documentclass[12pt]{article}

\usepackage[a4paper, margin=1in]{geometry}
\usepackage{xcolor}
\usepackage{tcolorbox}
\usepackage{fontawesome5} % icons
\usepackage{hyperref}
\usepackage{graphicx}
\usepackage{comment}
\usepackage[nameinlink,capitalize,noabbrev]{cleveref}
\usepackage{subcaption} % modern replacement for the deprecated subfigure package

\setlength{\parindent}{0pt}

% --- COLOR PALETTE (customize these) ---
\definecolor{maincolor}{HTML}{2A9D8F} % your teal
\definecolor{accentcolor}{HTML}{E76F51}
\definecolor{lightcolor}{HTML}{FFDEAD}

% Tell cleveref how to name your counter
\crefname{exercise}{assignment}{assignments}
\Crefname{exercise}{Assignment}{Assignments}

% --- ASSIGNMENT STYLE ---
\newcounter{exercise}

% Optional label: \exercisetitle[<label>]{<Title>}
\newcommand{\exercisetitle}[2][]{%
  \refstepcounter{exercise}% make this step referenceable/hyperlinked
  \vspace{0.5cm}%
  \colorbox{maincolor}{%
    \makebox[\dimexpr\linewidth-2\fboxsep][l]{%
      \textcolor{white}{\faPen\ Assignment \theexercise:~ #2}%
    }%
  }%
  % If an optional label was provided, set it
  \if\relax\detokenize{#1}\relax\else\label{#1}\fi
  \vspace{0.3cm}%
}


% --- IMPORTANT BOX ---
\newtcolorbox{importantbox}{
  colback=lightcolor,
  colframe=accentcolor,
  arc=4mm,
  boxrule=1pt,
  left=6pt,
  right=6pt,
  top=6pt,
  bottom=6pt
}

% --- PROMPT BOX ---
\newtcolorbox{promptbox}{
  colback=white,
  colframe=maincolor,
  arc=4mm,
  boxrule=1pt,
  left=6pt,
  right=6pt,
  top=6pt,
  bottom=6pt
}

% --- QUESTION BOX (with starred variant and cleveref names) ---
% Requires: \usepackage{hyperref} for \Form/\TextField
%           \usepackage{fontawesome5} for \faEdit

% Cleveref names for the question counter
\crefname{question}{question}{questions}
\Crefname{question}{Question}{Questions}

\makeatletter
\newcounter{question}
\newlength\questionheight
\setlength\questionheight{5cm} % default answer box height

% \question{<text>}                  -> with TextField (default height)
% \question[<height>]{<text>}        -> with TextField (custom height)
% \question*{<text>}                 -> no TextField, still numbered
\newcommand{\question}{\@ifstar{\question@text}{\question@field}}

% Common heading (increments counter and creates ref anchor)
\newcommand{\question@heading}[1]{%
  \refstepcounter{question}%
  \vspace{0.3cm}%
  \noindent\textbf{\faEdit\ Question \thequestion: #1}%\\[0.3em]%
}

% Starred form: numbered question without input field
\newcommand{\question@text}[1]{%
  \question@heading{#1}%
%  \vspace{0.3cm}%
}

% Unstarred form: numbered question with input field (optional height)
\newcommand{\question@field}[2][\questionheight]{%
  \question@heading{#2}%
 \\[0.3em] \begin{Form}
    \TextField[
      name=Q\thequestion,
      multiline=true,
      width=\linewidth,
      height=#1,
      charsize=10pt,
      bordercolor={0 0 0},
      borderwidth=1.8,
      backgroundcolor={1 1 1}
    ]{}
  \end{Form}
  \vspace{0.3cm}%
}
\makeatother

% --- SOLUTION BOX ---
%--- Definition of solution/no solution env to change between the teacher and the student versions --------%

% --- Teacher toggle ---
\newif\ifteacher
\teachertrue  % Write teacherfalse for student version and teachertrue for teachers version

% --- Solution environment ---
\newtcolorbox{solutionbox}{colback=gray!10,colframe=black!60!black,arc=1mm,boxrule=1pt}

\ifteacher
\newenvironment{solution}{%
  \begin{solutionbox}\textbf{Solution:}
}{%
  \end{solutionbox}
}
\else
  \excludecomment{solution}
\fi



\title{\color{maincolor} Style Card}
\date{\vspace{-5ex}}
\begin{document}

\maketitle
\vspace{-5ex}
\subsection*{Initial impression}
\question[2\baselineskip]{Group members' names (2--3)}

\question[1\baselineskip]{When do you guess that this artwork was created?}

\question[3\baselineskip]{Do you have an idea what art style it could be?}

\question[3\baselineskip]{What do you like or dislike about the artwork}

\question{Using the Style Guide, write down three style features of this artwork}


\subsection*{Art Style identification}
For the upcoming questions about the identification of the style, you can use an LLM assistent to help you. 


Some helpful prompts:
\begin{promptbox}
“Based on this description of a painting: [1–2 sentences describing what you see], identify likely art style(s). List 3 signature visual features, 2 artists connected to that style, and a typical date range. Include 1 museum link to verify. Keep under 100 words.”
\end{promptbox}
\begin{promptbox}
“List 3 public-domain artworks that exemplify [style], with museum links.”
\end{promptbox}
\begin{promptbox}
“Provide 2 museum or academic links that confirm these features for [style]. Summarize what each source confirms in one sentence.”
\end{promptbox}

\begin{importantbox}
    Make sure that you verify your information with trustworthy sources, for example museum web pages.
\end{importantbox}


\question[3\baselineskip]{What are the artist and title of the artwork?}

\question[3\baselineskip]{What art style(s) does the artwork belong to?}

\question[2\baselineskip]{What date range did this art style thrive?}

\question[2\baselineskip]{What region or movement does this art style originate from?}

\question[6\baselineskip]{List three visual hallmarks of the style that can be seen in your mystery image}

\question[4\baselineskip]{Two other artists associated with this style}

\question[4\baselineskip]{One other representative artwork (not the mystery artwork) (title + artist + date)}

\question[6\baselineskip]{Why did this style emerge? (e.g., social change, technology, philosophy)}

\subsection*{Citations (paste URLs)}
\question[3\baselineskip]{Source 1 (museum/academic)}
\question[3\baselineskip]{Source 2 (optional)}

\end{document}